%% Incipit %%

Le 30 mai 2023, Gallica fêtait son dix millionième document numérisé\footcite[]{gallica10}. Exposant par la même occasion le succès de la numérisation des documents patrimoniaux, alors que 23 ans auparavant, elle peinait toujours à trouver sa place dans des institutions de conservation comme la BNF. Aujourd'hui cette numérisation est devenue indissociable de la communication des ces documents au plus grand nombre. Pourtant, une des conséquences de cette numérisation en masse est l'effacement de certains de ces documents parmi ces quantités de documents massive. En l'absence d'une représentation physique, il revient à ces institutions de mettre en place les outils nécessaires à la réalisation de leur mission de communication des collections patrimoniales dont ils ont la responsabilité. Ils disposent pour cela de nombreuses années d'expérience à documenter leur collection dans un objectif de communication. Malheureusement pour eux, cette numérisation est seulement un des témoins des bouleversements provoqués par l'informatisation de notre société. Les attentes des publics ont changé et une adaptation est nécessaire. 

%% Introduction INA %%

C'est dans ce contexte propices aux réflexions sur la communication de ces documents que prend place notre stage à l'Institut National de l'audiovisuel. Né de la dissolution de l’ORTF en 1974, l’Institut National de l’Audiovisuel, abrégé INA, est aujourd’hui une institution majeure de la conservation du patrimoine français. Pourtant, à sa naissance, il réunissait les oubliés de la dissolution avec l’archivage, la formation, la production de création et la recherche\footcite[][3]{hoog2006}. Depuis sa création, l’institut a conservé toutes ses missions d’origine mais s’est vu transformé par plus de 50 ans de mutation des institutions patrimoniales ainsi que de la définition même de la notion de « patrimoine », auquel l’audiovisuel s’est progressivement greffé. Porte étendard de ces évolutions à l’INA, la loi du dépôt légal du 20 juillet 1992 grâce à laquelle l’institution intègre deux nouvelles missions : la conservation patrimoniale et la communication à des fins scientifiques\footcite{fillioud1994}. D’un autre côté, en matière de changements, le paysage télévisuel n’est pas en reste non plus. Les chaînes se multiplient, les contenus se diversifient, l’INA a adapté son fonctionnement à ces évolutions. A titre général, le monde s’est progressivement informatisé, ce qui permis au grand public de se rapprocher de ses archives audiovisuelles grâce à des plateformes comme ina.fr\footcite{zotero-item-158} alors que les professionnels de l’INA, eux, perdent le contact avec les supports analogiques pour l’immatérialité de la donnée\footcite[][68]{bassaïstéguy2025}. Aujourd’hui professionnel de l’INA ou grand public, tous accèdent aux documents par les données qui leur sont associées\footcite[][118]{alquier2017}. En cela le \gls{td} réalisé de ces documents tient une place importante dans leur communication à tous les publics. Des années durant, les politiques documentaires ont adapté leurs pratiques à ses usagers. La conséquence étant aujourd'hui que cette quantité massive de données est empreinte de tous les bouleversements d'ordre sociaux, économique ou technologiques au cours desquelles elle se construisait. A tel point où désormais, l'accessibilité aux collections des institutions patrimoniales est un véritable enjeu. 

%% Stage %%

C'est de ce constat que né le projet auquel j'ai pris part au cours de mon stage à l'institut national de l'audiovisuel. J'avais pour mission de participer aux réflexions sur la conceptualisation d'un outil permettant de rendre ces données moins opaques, plus loquaces afin de faciliter la découvrabilité des fonds et la recherche dans les bases de données de l'INA. Le \gls{td} de ces documents audiovisuels n'est pas égal, il est le résultat de multiples politiques documentaires, découlant d'un contexte spécifique, sans pour autant avoir laissé de trace de ces changements. En conséquent, la précision de ce \gls{td} varie selon la \gls{pd} en vigueur lors de son dernier traitement mais aussi selon un certain nombre de critères variables mais qui ne sont aucunement documenté non plus. C'est pourquoi ce projet vise à cartographier l'histoire du \gls{td} à l'INA. Parmi toutes les données que produit l'INA, ce sont à celles-ci que sont confronté les chercheurs et ce sont celles-ci qui déterminent la compréhension des collections de l'INA. Leur absence est aussi équivoque que leur présence. Ce projet ayant pour objectif de produire un outil permettant à différents publics de comprendre, explorer et exploiter les bases de données de l'INA, il nous a été essentiel de déterminer \textbf{en quoi l'évolution des pratiques documentaires au sein des institutions patrimoniales représente aujourd'hui un enjeu d'accessibilité aux collections ? Quels outils permettent d'adapter cet héritage à de nouvelles habitudes de consultation ?}

%% Plan %% 

Afin de répondre à cette question, nous procéderons en trois étapes. Premièrement nous analyserons les publics confrontés à ces collections patrimoniales et nous expliciterons à quels enjeux répond leur accessibilité. Une fois que nous aurons déterminer cela, nous parcourrons l'histoire du \gls{td} à partir de la naissance de cette institution afin de déterminer comment le \gls{td} en est arrivé à générer la situation précédemment décrite. Pour finir, nous utiliserons ce que nous aurons déduit de ces analyses afin de déterminer quels outils nous permettent d'adapter cette histoire du \gls{td} aux nouvelles habitudes de consultation de ces publics. 

%% Méthodologie %%

Cette analyse utilisera principalement des exemples issus de la situation actuelle du \gls{td} à l'institut national de l'audiovisuel desquels nous tenterons de tirer des conclusions générales sur des institutions patrimoniales issues du paysage patrimonial français, en prêtant une attention particulière aux institutions responsables de la mission du dépôt légal. Nous tirerons nos conclusions de l'analyse de la littérature scientifique, de notre expérience personnelle au sein de l'institution ainsi que de l’expérience de trois professionnels de l'INA exprimées dans trois entretiens semi-dirigés. Ils ont été sélectionnés afin de correspondre à trois profils de publics ciblés par les activités de l'INA ou concernés par le \gls{td} : les documentalistes avec une documentaliste ayant souhaité gardé l'anonymat, du service \enquote{Traitement et enrichissement des collections}, les chercheurs avec Arthur Hénaff, chef de service de l'INAthèque et enfin les professionnels de l'audiovisuel avec Pascal Lebeurier, cadre documentaire du service \enquote{Clients et partenaires}. 

%% Vocabulaire %%

Au cours de cet exposé nous utiliserons à plusieurs reprises le terme d' \enquote{archives} pour parler de documents audiovisuels conservés. Ce terme n'est pas utilisé au sens premier des archives comme le résultat d'une activité administrative ou juridique. Exit également la notion d'original, incompatible avec ces documents, majoritairement numériques pour lesquels la question d'original est très complexe. Nous utiliserons ici le terme d'\enquote{archives} comme un objet conservé et communiqué par une institution\footcite[Réflexion inspirée de l'article suivant][]{zotero-item-739}. 